\selectlanguage{ngerman}%
\frenchspacing%
\begin{ledgroupsized}[r]{120mm}%
\footnotesize%
\pstart%
\noindent\textbf{Überlieferung:}%
\pend%
\end{ledgroupsized}%
\begin{ledgroupsized}[r]{114mm}%
\footnotesize%
\pstart%
\parindent -6mm%
\makebox[6mm][l]{\textit{E\textsuperscript{1}}}%
Druck:
\textit{AL2%
}, centuriae III--IV (1715), Anhang, Nr.\;3, S.\;49--51.\cite{01457}\cite{01458}
\pend%
\pstart%
\parindent -6mm%
\makebox[6mm][l]{\textit{E\textsuperscript{2}}}%
Druck (nach \textit{E\textsuperscript{1}}):
\textit{LOD} II.2, S.\:78--79.\cite{00150}
\pend%
\pstart%
\parindent -6mm%
\makebox[6mm][l]{\textit{Ü}}%
(ins Französische):
\textsc{Peyroux} 1990, S.\:110--111.\cite{01414}
\pend%
\end{ledgroupsized}%
\selectlanguage{latin}%
\frenchspacing%
\vspace{8mm}%
\normalsize%
\pstart%
\noindent%
\lbrack S.\:49\rbrack%
\hfill%
lllustris Dn.\;G.\;G.\;Leibnitii%
\indexN{\textls{Leibniz} (Leibnitius, GGL), Gottfried Wilhelm 1646\textendash1716}
\hfill%
\phantom{\lbrack S.\:49\rbrack}%
\pend%
\vspace*{0.5em}%
\pstart%
\centering%
De Aetiologia Barometri,%
\indexS{aetiologia barometri}%
\indexS{barometrum}\\
\edtext{%
occasione n.\;3 et 4 appendicis ad Centur.\;1.\;et 2.\\
Ephemeridum novarum Societatis Caesareae Naturae Curiosorum%
\indexS{Academia naturae curiosorum}%
}{%
\lemma{occasione \lbrack...\rbrack\ Curiosorum}%
\Cfootnote{%
Vgl.
\textit{AL2}, centuriae I--II (1712), Anhang,
Nr.\;3, S.\;92--98\cite{01457}
(B.\;\textsc{Ramazzini}, \glqq Epistola de causa adscensus et descensus mercurii in fistula Torricelliana\grqq);\cite{01460}
Nr.\;4, S.\;98--110\cite{01457}
(G.\:C.\;\textsc{Schel\-ham\-mer}, \glqq Notata ad hanc epistolam\grqq).\cite{01461}
Ramazzinis \glqq Epistola\grqq\ ist ein Nachdruck seines Briefes an L.\;Schröck vom 1.\;Juni 1710,
nach B.\;\textsc{Ramazzini}, \textit{Ephemerides barometricae. Cum disquisitione}
(Padua 1710, S.\;191--202).\cite{01462}%
}}
\pend%
\vspace*{1.0em}%
\pstart%
\noindent%
Barometrum,%
\indexS{barometrum}
quod vocant,
Hydrargyro%
\indexS{hydrargyrum in barometro}
aeris pondus indicans,%
\indexS{pondus aeris}
\edtext{%
primus invenit Torricellius,%
\indexN{\textls{Torricelli} (Torricellius), Evangelista 1608\textendash1647}%
}{%
\lemma{primus invenit Torricellius}%
\Cfootnote{%
Siehe \textsc{E.\;Torricelli},
Brief an M.\;Ricci vom 11.\;Juni 1644\cite{00107}
(in C.\,\textsc{Dati},
\textit{Lettera a Filaleti ... della famosissima esperienza dell’argento vivo},
Florenz 1663, S.\;20\,f.;\cite{00032}
\textit{TO} III, S.\;186--188).\cite{00108}%
}}
\edtext{%
primus Gerikius%
\indexN{\textls{Guericke} (Gerickius, Gericke), Otto von 1602\textendash1686}
sub nomine virunculi meteorologi
aliquandiu occultatum%
\indexS{virunculus meteorologus}
ad praedictiones mutationum aeris%
\indexS{praedictio mutationum aeris}%
\indexS{mutatio aeris}
publice adhibuit.%
}{%
\lemma{primus \lbrack...\rbrack\ adhibuit}%
\Cfootnote{%
Siehe O.\;v.\;\textsc{Guericke}, \textit{Experimenta nova},%
\cite{00055}
lib.\;III, cap.\;20; 31; lib.\;V, appendix
(Amsterdam 1672, S.\;100a--b; 115a--b; 188b--189a; 195b--196a).
Leibniz hatte 1672 diese Stellen exzerpiert:
Vgl. \textit{LSB} VIII,\;1 N.\;36,%
\cite{01197}
S.\;257.23--30; 259.1--10; 270.14--25.%
}}
\edtext{%
Sed cum observaretur
aerem,%
\indexS{aer factus levior}
imminente pluvia,%
\indexS{pluvia imminens}
leviorem fieri solere%
\lbrack,\rbrack\
siccum vero et serenum,%
\indexS{aer siccus}%
\indexS{aer serenus}
ponderosiorem ostendi;%
\indexS{aer ponderosior}
cum contra tempore ad pluviam vergente,%
\indexS{tempus vergens ad pluviam}%
\indexS{pluvia imminens}
vaporibus aqueis%
\indexS{vapor aqueus}
magis oneratus%
\indexS{aer oneratus vaporibus aqueis}
videri posset;
problema perplexum%
\indexS{problema perplexum}
varie agitatum est,%
}{%
\lemma{Sed cum \lbrack...\rbrack\ agitatum est}%
\Cfootnote{%
Auf die vorherige wissenschaftliche Diskussion über die wetterbedingten Änderungen des Luftgewichts geht ausführlich
-- in Anlehnung an B.\;\textsc{Ramazzini}, \textit{Ephemerides barometricae Mutinenses%
} (Modena 1695) --\cite{01516}
die Aufzeichnung N.\;\ref{60267} ein.%
}}
et in primis egregii Medici%
\indexS{medicus}
et Naturae consulti,%
\indexS{vir consultus naturae}%
\indexS{natura}
mihique amici,%
\indexS{amicus}
\edtext{%
Ramazzinus%
\indexN{\textls{Ramazzini} (Ramazzinus), Bernardino 1633\textendash1714}
et Schelhamerus%
\indexN{\textls{Schelhammer} (Schelhammerus), Günther Christoph 1649\textendash1716}
de ea re inter se contulerunt.%
}{%
\lemma{Ramazzinus \lbrack...\rbrack\ contulerunt}%
\Cfootnote{%
Anspielung auf die Auseinandersetzung zwischen B.\;Ramazzini und G.\:C.\;Schelhammer,
welche ab 1697 in zahlreichen Veröffentlichungen %
ihren Niederschlag fand.
Siehe für die Einzelheiten die editorische Einleitung zu N.\;\ref{paradoxon_intro} (\refpassage{paradoxon_RamaSchel-fhg-1}{paradoxon_RamaSchel-fhg-2}).%
}}
Mihi in mentem jam olim venerat,
pluvia imminente%
\indexS{pluvia imminens}
levari aeris columnam,%
\indexS{columna aeris}
quod ea guttulas%
\indexS{guttula cadens}
jam jam cadentes %
et pluviae mox materiam%
\indexS{materia pluviae}%
\indexS{pluvia imminens}
praebituras %
non amplius sustineret;
contra sereno tempore%
\indexS{tempus serenum} %
aquam%
\indexS{aqua dilatata in vapores}%
\indexS{aqua mixta aeri}
in minutissimos vapores%
\indexS{vapor aqueus}
dilatatam%
\indexS{aqua dilatata in vapores} %
misceri aeri,%
\indexS{aqua mixta aeri}
nec descendendi%
\indexS{vis descendendi}
aeremve penetrandi vim%
\indexS{vis penetrandi aerem}
habere%
\lbrack,\rbrack\
ac proinde una cum columna%
\indexS{columna aeris}
a subjecto Mercurio%
\indexS{mercurius subjectus columnae aeris}
sustentari.
\edtext{%
Proposui et experimentum,%
\indexS{experimentum propositum}
quo ad oculum%
\indexS{oculus}
appareret %
gravia solida%
\indexS{grave solidum}%
\indexS{grave descendens in fluido}
descendentia in gravi fluido%
\indexS{fluidum grave}
non gravitare in fundum fluido%
\indexS{fluidum grave}
subjectum.%
}{%
\lemma{Proposui \lbrack...\rbrack\ subjectum}%
\Cfootnote{%
Das Gedankenexperiment wird in den Aufzeichnungen N.\;\ref{53442}, \ref{60267} und \ref{36805} dargelegt.
Siehe zudem G.\:W.\;\textsc{Leibniz}, Brief an F.\;Hoffmann vom 19. März 1701
(\textit{LSB} III,\;8 N.\;225, S.\;577.17\,f.).\cite{01456}%
}}
Nempe tubus aqua plenus%
\indexS{tubus plenus aqua}%
\indexS{aqua in tubo}
in una parte librae bilancis%
\indexS{libra bilanx}%
\indexS{bilanx}
suspenditur,%
\indexS{tubus suspensus in libra}
et gravi aliquo aquae innatante%
\indexS{grave innatans aquae}
aut aliter in ea suspenso%
\indexS{grave suspensum in aqua}
et tubum%
\indexS{tubus plenus aqua}
vel lancem%
\indexS{lanx librae}
gravante,
totum per pondus%
\indexS{pondus redactum ad aequilibrium}
in opposita lance%
\indexS{lanx librae}
ad aequilibrium%
\indexS{aequilibrium in libra}
redigitur:
mox eodem gravi in aqua descendente%
\indexS{grave descendens in aqua}
(sive resecto filo%
\indexS{filum resectum}
sive aqua in cavitatem gravis admissa)%
\indexS{aqua admissa in cavitatem}%
\indexS{cavitas repleta aqua}
levior fit lanx%
\indexS{lanx librae}
tubi,%
\indexS{tubus suspensus in libra}
et pondus oppositum%
\indexS{pondus oppositum in libra}
praeponderat,%
\indexS{pondus praeponderans}
donec grave cadens%
\indexS{grave cadens}
fundum at%
\lbrack S.\:50\rbrack%
tingat.
Itaque occasione data
\edtext{%
sententiam meam%
\indexS{sententia}
Celeberrimo Ramazzino%
\indexN{\textls{Ramazzini} (Ramazzinus), Bernardino 1633\textendash1714}
per literas%
\indexS{literae}
significavi.%
}{%
\lemma{sententiam \lbrack...\rbrack\ significavi}%
\Cfootnote{%
Siehe G.\:W.\;\textsc{Leibniz}, Brief an B.\;Ramazzini vom 18.\;März 1700
(\textit{LSB} III,\;8 N.\;139, S.\;371\,f.).%
}}
Virum acutum%
\indexS{vir acutus}
et rationi%
\indexS{ratio}
cedere facilem
tetigit meditatio mea,%
\indexS{meditatio}
\edtext{%
cumque etiam experimentum,%
\indexS{experimentum Ramazzini}
adjuvante
\edtext{%
Cl.\;Gratiano,%
\indexN{\textls{Graziani} (Gratianus), Giovanni, gest. 1744}%
}{%
\lemma{Cl.\;Gratiano}%
\Cfootnote{%
Giovanni Graziani lehrte seit 1703 in Padua\indexO{Padua} (zunächst Astronomie, ab 1709 Philosophie).%
}}
sumsisset,
laudabili candore
publice professus est %
veritati%
\indexS{veritas}
a se manus dari.%
\indexS{manus data veritati}%
}{%
\lemma{cumque \lbrack...\rbrack\ dari}%
\Cfootnote{%
Siehe \textsc{Ramazzini}, Brief an Schröck vom 1.\;Juni 1710,\cite{01460}
in
\textsc{Ders.}, \textit{Ephemerides barometricae. Cum disquisitione}
(1710, S.\;197--201)\cite{01462}
sowie in
\textit{AL2}%
, 1712 (Anhang, %
S.\;95--97).\cite{01457}
Das Experiment, auf das Leibniz anspielt, hatte nach Ramazzinis Angabe zur Fastnachtszeit 1710 in Padua stattgefunden.
Ramazzini lehrte seit 1700 Medizin an der dortigen Universität.%
\indexO{Padua}%
}}
Sed
\edtext{%
Antagonistae%
\indexN{\textls{Schelhammer} (Schelhammerus), Günther Christoph 1649\textendash1716}%
\indexS{antagonista}%
}{%
\lemma{Antagonistae}%
\Cfootnote{%
Gemeint ist Günther Christoph Schelhammer.%
}}
haud minus celebri nondum satisfactum video.
\pend%
\pstart%
\edtext{%
Objicit%
\indexS{objectio}
\textls{primo,}
vapores aquosos%
\indexS{vapor aquosus}
natantes in aere%
\indexS{vapor natans in aere}
eum premere;%
\indexS{vapor premens aerem}
at corpus suspensum ex filo%
\indexS{corpus suspensum e filo}
in experimento D.\;Ramazzini%
\indexN{\textls{Ramazzini} (Ramazzinus), Bernardino 1633\textendash1714}%
\indexS{experimentum Ramazzini}
non premere aquam.%
\indexS{aqua in tubo}%
}{%
\lemma{Objicit \lbrack...\rbrack\ aquam}%
\Cfootnote{%
Vgl. \textsc{Schelhammers} kritische Bemerkungen
in \textit{AL2}, %
1712 (Anhang, %
S.\;98).\cite{01461}\cite{01457}%
}}
Sed responsio est facilis%
\lbrack:\rbrack%
\indexS{responsio objectioni}
discrimen%
\indexS{discrimen}
inde nullum oriri ad rem faciens%
\lbrack,\rbrack\
sufficit enim
quod utrobique corpus grave in fluido descendens%
\indexS{corpus descendens in fluido}%
\indexS{grave descendens in fluido}%
\indexS{fluidum grave}
cum eo non gravitet;%
\indexS{corpus non gravitans}
quod tamen antea gravitabat in lancem%
\indexS{corpus gravitans in lancem}%
\indexS{lanx librae}
oppositumque pondus%
\indexS{pondus oppositum in libra}
sive a filo sive ab aqua sustentaretur.%
\indexS{grave sustentum a filo}%
\indexS{corpus sustentum a filo}%
\indexS{grave sustentum ab aqua}%
\indexS{corpus sustentum ab aqua}
\edtext{%
Objicit%
\indexS{objectio}
\textls{secundo,}
hac causa supposita%
\indexS{causa supposita}
Mercurium ergo non descensurum,%
\indexS{mercurius descendens in barometro}
nisi cum actu%
\indexS{actus}
pluit%
\indexS{pluvia actualis}%
\lbrack.%
}{%
\lemma{Objicit \lbrack...\rbrack\ pluit}%
\Cfootnote{%
Vgl. a.a.O., S.\;98\,f.\cite{01461}\cite{01457}%
}}
S\rbrack ed respondetur%
\indexS{responsio objectioni}%
\lbrack:\rbrack\
etsi aqueae particulae%
\indexS{particulae aqueae separatae ab aere}
ab aere separari
et cadere incipiant,%
\indexS{particulae aqueae cadentes}
non tamen statim ad nos pervenire
aut pluviam formare;%
\indexS{particulae aqueae formantes pluviam}%
\indexS{pluvia formata particulis aqueis}
sed interdum a vento%
\indexS{ventus agitans aerem}%
\indexS{ventus dissipans particulas aqueas}%
\indexS{ventus retardans particulas aqueas}
in superiore regione%
\indexS{regio aeris superior}
aerem agitante,%
\indexS{aer regionis superioris}%
\indexS{aer agitatus a vento}
vel alia causa,%
\indexS{causa alia}
dissipari%
\indexS{particulae aqueae dissipatae a vento}
aut rerardari.%
\indexS{particulae aqueae retardatae a vento}
\edtext{%
Objicit%
\indexS{objectio}
\textls{tertio,}
Mercurium saepe descendere%
\indexS{mercurius descendens in barometro}
et varie mutari,
nulla pluvia cadente.%
\indexS{pluvia non cadens}%
}{%
\lemma{Objicit \lbrack...\rbrack\ cadente}%
\Cfootnote{%
Vgl. a.a.O., S.\;99.\cite{01461}\cite{01457}%
}}
Solutio%
\indexS{solutio obiectionis}
ex praecedenti patet:
nam fieri potest et solet,
ut particulae descendere incipientes%
\indexS{particulae aqueae descendentes in aere}
rursus sole%
\indexS{sol dissipans particulas aqueas}
vel vento,%
\indexS{ventus dissipans particulas aqueas}
vel alia causa,%
\indexS{causa alia}
dissipentur,%
\indexS{particulae aqueae dissipatae}
imo plane alio ferantur,
novo aere substituto;%
\indexS{aer novus substitutus}
\edtext{%
quod etiam jam notavit Cl.\;Ramazzinus.%
\indexN{\textls{Ramazzini} (Ramazzinus), Bernardino 1633\textendash1714}%
}{%
\lemma{quod \lbrack...\rbrack\ Ramazzinus}%
\Cfootnote{%
Siehe \textsc{Ramazzini}, Brief an Schröck vom 1.\;Juni 1710\cite{01460}
in \textsc{Ders.}, \textit{Ephemerides barometricae. Cum disquisitione}
(1710, S.\;200\,f.)\cite{01462}
sowie in \textit{AL2}%
, 1712 (Anhang, %
S.\;97).\cite{01457}%
}}
\edtext{%
Objicit%
\indexS{objectio}
\textls{quarto,}
alias magnas mutationes in Barometro%
\indexS{mutatio in barometro}%
\indexS{barometrum}
notari ex aliis causis,%
\indexS{causa alia}%
\indexS{causa mutationis in barometro}
verbi gratia %
a validis procellis.%
\indexS{procella}%
}{%
\lemma{Objicit \lbrack...\rbrack\ procellis}%
\Cfootnote{%
Vgl. \textsc{Schelhammers} kritische Anmerkungen\cite{01461}
in \textit{AL2}%
, 1712 (Anhang, %
S.\;99).\cite{01457}%
}}
Hoc concedo,
neque enim omnes mutationes%
\indexS{mutatio ponderis columnae aereae}
ponderis%
\indexS{pondus columnae aereae}
columnae aereae%
\indexS{columna aerea}
ex hac sola,
quam adduxi,%
\indexS{causa adducta}
causa%
\indexS{causa mutationis columnae aereae}
oriri puto;
sufficit eam saepe valere%
\indexS{causa valens saepe}
et principali phaenomeno%
\indexS{phaenomenon principale}
\edtext{%
(ut jam Dn.\;Ramazzinus ait)%
\indexN{\textls{Ramazzini} (Ramazzinus), Bernardino 1633\textendash1714}%
}{%
\lemma{ut \lbrack...\rbrack\ ait}%
\Cfootnote{%
Siehe \textsc{Ramazzini}, Brief an Schröck vom 1.\;Juni 1710,\cite{01460}
in \textsc{Ders.}, \textit{Ephemerides barometricae. Cum disquisitione}
(1710, S.\;201)\cite{01462}
sowie in \textit{AL2}%
, 1712 (Anhang, %
S.\;97\,f.).\cite{01457}%
}}
satisfacere,
cum scilicet Mercurii%
\indexS{mercurius descendens in barometro}%
\indexS{mercurius ascendens in barometro}
descensus%
\indexS{descensus mercurii in barometro}
pluviam%
\indexS{descensus mercurii nuntians pluviam}%
\indexS{pluvia nuntiata a barometro}
aut nivem,%
\indexS{descensus mercurii nuntians nivem}%
\indexS{nivis nuntiata a barometro}
ascensus%
\indexS{ascensus mercurii in barometro}
serenitatem nuntiat.%
\indexS{ascensus mercurii nuntians serenitatem}%
\indexS{serenitas nuntiata a barometro}
Sane ventus vehemens%
\indexS{ventus suspendens pondus aeris}
etiam aeris pondus suspendere potest,%
\indexS{pondus aeris suspensum a vento}
unde jam
\edtext{%
Gerikius%
\indexN{\textls{Guericke} (Gerickius, Gericke), Otto von 1602\textendash1686}
notavit,
in magnis tempestatibus%
\indexS{tempestas magna}
Mercurium Barometri descendere.%
\indexS{mercurius descendens in barometro}%
\indexS{descensus mercurii in barometro}%
\indexS{barometrum}%
}{%
\lemma{Gerikius \lbrack...\rbrack\ descendere}%
\Cfootnote{%
Siehe \textsc{Guericke}, \textit{Experimenta nova},%
\cite{00055}
lib.\;III, cap.\;20;
lib.\;V, appendix
(100b; 189a; 195b--196a).
Vgl. zudem Leibnizens Exzerpt,
\textit{LSB} VIII,\;1 N.\;36,%
\cite{01197}
S.\;270.17--19.%
}}
Sed et aliter mutare venti%
\indexS{ventus mutans pondus aeris}
pondus columnae aereae%
\indexS{pondus columnae aereae}%
\indexS{columna aerea}%
\indexS{pondus aeris mutatum a vento}
possunt,
ex.\:gr.\;duo venti sibi occurrentes%
\indexS{venti occurrentes sibi}%
\indexS{venti comprimentes aerem}
aerem interceptum%
\indexS{aer interceptus}
\lbrack S.\:51\rbrack\
compriment%
\indexS{aer compressus a ventis}
et reddent densiorem%
\indexS{aer factus densior}
adeoque graviorem;%
\indexS{aer factus gravior}
contra in diversa tendentes,%
\indexS{venti tendentes in diversa}
distrahent eundem%
\indexS{venti distrahentes aerem}%
\indexS{aer distractus a ventis}
et reddent rariorem%
\indexS{aer factus rarior}
et leviorem;%
\indexS{aer factus levior}
aer autem rarior guttas aqueas%
\indexS{gutta aquea}
minus sustentabit,%
\indexS{aer sustentans guttas aqueas}
et ita ex duplici capite
levior erit.%
\indexS{aer factus levior}
Harum aliarumque causarum%
\indexS{causa alia}
interventio%
\indexS{interventio aliarum causarum}
facit,
ut regulae praedictionum ex Barometro%
\indexS{regula praedictionum e barometro}%
\indexS{praedictio e barometro}%
\indexS{barometrum}
multas exceptiones patiantur.%
\indexS{exceptio regulae}
\edtext{%
Objicit%
\indexS{objectio}
tandem \textls{quinto,}
Experimentum Ramazzinianum%
\indexN{\textls{Ramazzini} (Ramazzinus), Bernardino 1633\textendash1714}%
\indexS{experimentum Ramazzini}
nonnihil differre a meo.
Nam ego corpus cavum proposueram%
\indexS{corpus cavum natans}
quod nataret in aqua Tubi,%
\indexS{tubus plenus aqua}%
\indexS{aqua in tubo}
mox aqua repletum%
\indexS{corpus cavum repletum aqua}
in ea mergeretur;
Ramazzinus%
\indexN{\textls{Ramazzini} (Ramazzinus), Bernardino 1633\textendash1714}
suum grave ex filo suspendit,%
\indexS{grave suspensum e filo}
et mox filo resecto%
\indexS{filum resectum}
cadere permisit.%
\indexS{grave cadens}%
}{%
\lemma{Objicit \lbrack...\rbrack\ permisit}%
\Cfootnote{%
Vgl. \textsc{Schelhammers} kritische Anmerkungen\cite{01461}
in \textit{AL2}%
, 1712 (Anhang, %
S.\;99).\cite{01457}%
}}
Sed respondeo
primo
nullam esse rationem differentiae,%
\indexS{ratio differentiae}
ut ostensum est ad objectionem primam;%
\indexS{objectio}
secundo
si Cl.\;objector%
\indexN{\textls{Schelhammer} (Schelhammerus), Günther Christoph 1649\textendash1716}%
\indexS{objector}
sumere experimentum dignetur,%
\indexS{experimentum}
cavo natante debite substituto%
\indexS{corpus cavum natans}
in locum gravis suspensi ex filo,%
\indexS{grave suspensum e filo}
eosdem eventus reperturum,%
\indexS{eventus}
cum eadem ratio subsit.
\pend%
\pstart%
Obiter addo
\edtext{%
in Epistolae%
\indexS{epistola}
meae secunda linea%
\indexS{linea textus}
(p.\;95)
vocem%
\indexS{vox delenda}
\textit{mehercle} esse delendam,%
}{%
\lemma{in Epistolae \lbrack...\rbrack\ delendam}%
\Cfootnote{%
Ramazzini hat in seinem Brief an Schröck vom 1.\;Juni 1710\cite{01460}
auch den Brief abgedruckt, mit dem ihm Leibniz am 18. März 1700 seine Erklärung der wetterbedingten Änderung des Luftgewichts gesendet hatte (vgl. \textit{LSB} III,\;8, N.\;139, S.\;371\,f.).\cite{01454}
In dem von Ramazzini veranlassten Abdruck dieses Briefes findet sich an der von Leibniz angegebenen Stelle tatsächlich der Fehler, um dessen Berichtigung er in N.\;\ref{54958} bittet;
vgl.
\textsc{Ramazzini}, \textit{Ephemerides barometricae. Cum disquisitione}
(1710, S.\;196)\cite{01462}
sowie \textit{AL2}%
, 1712 (Anhang, %
S.\;95).\cite{01457}
Leibniz zitiert offenbar aus den \textit{Ephemerides} der Leopoldina.%
}}
ex alia credo non bene scripta lectaque natam,
\edtext{%
quae qualis fuerit nunc non divino.%
}{%
\lemma{quae \lbrack...\rbrack\ divino}%
\Cfootnote{%
Anstelle von \textit{mehercle meis} stand in Leibnizens Brief an Ramazzini vom 18. März 1700 das Wort \textit{Mechanicis};
vgl. \textit{LSB} III,\;8 N.\;139, S.\;371.13.\cite{01454}%
}}
\pend%
\selectlanguage{latin}%
\frenchspacing%
\normalsize%
